\documentclass[paper=a4, fontsize=11pt]{article}

%----------------------------------------------------------------------------------------
%	PACKAGES AND OTHER DOCUMENT CONFIGURATIONS
%----------------------------------------------------------------------------------------

\usepackage{amsmath,amsfonts,amsthm} % Math packages
\usepackage{sectsty} % Allows customizing section commands
\allsectionsfont{\centering \normalfont\scshape} % Make all sections centered, the default font and small caps

\usepackage{fancyhdr} % Custom headers and footers
\pagestyle{fancyplain} % Makes all pages in the document conform to the custom headers and footers
\fancyhead{} % No page header - if you want one, create it in the same way as the footers below
\fancyfoot[L]{} % Empty left footer
\fancyfoot[C]{} % Empty center footer
\fancyfoot[R]{\thepage} % Page numbering for right footer
\renewcommand{\headrulewidth}{0pt} % Remove header underlines
\renewcommand{\footrulewidth}{0pt} % Remove footer underlines
\setlength{\headheight}{13.6pt} % Customize the height of the header

\numberwithin{equation}{section} % Number equations within sections (i.e. 1.1, 1.2, 2.1, 2.2 instead of 1, 2, 3, 4)
\numberwithin{figure}{section} % Number figures within sections (i.e. 1.1, 1.2, 2.1, 2.2 instead of 1, 2, 3, 4)
\numberwithin{table}{section} % Number tables within sections (i.e. 1.1, 1.2, 2.1, 2.2 instead of 1, 2, 3, 4)

\setlength{\parindent}{4em}
\setlength{\parskip}{1em}
\renewcommand{\baselinestretch}{1.5}

\usepackage{graphicx}
\graphicspath{ {images/} }

\usepackage{xepersian}
\settextfont[Path=fonts/]{Vazir.ttf}
\setlatintextfont{Times New Roman}

%----------------------------------------------------------------------------------------
%	TITLE SECTION
%----------------------------------------------------------------------------------------

\newcommand{\horrule}[1]{\rule{\linewidth}{#1}} % Create horizontal rule command with 1 argument of height

\title{
\normalfont\normalsize
\includegraphics[scale=0.1]{aut}
\hspace{5cm}
\includegraphics[scale=0.1]{ceit} \\
\textsc دانشگاه صنعتی امیرکبیر \\
\textsc دانشکده مهندسی کامپیوتر و فناوری اطلاعات
\horrule{0.5pt} \\ [0.4cm] % Thin top horizontal rule
\huge بهینه‌سازی و کاربرد آن در شبکه‌های کامپیوتری \\
\huge تمرین سوم
\horrule{2pt} \\ [0.5cm] % Thick bottom horizontal rule
}

\author{پرهام الوانی}

\date{\normalsize\today} % Today's date or a custom date

\begin{document}

\maketitle % Print the title

\section{سوال اول}
\paragraph{}
از آنجایی که تابع هدف برای تمام گره‌ها یکسان است برای هر گره تنها محدودیت‌ها را می‌نویسیم.

\section{سوال دوم}
\subsection{الف}
\paragraph{}
متغیر $x_{i,(u,v)}$
نشان می‌دهد که یال $u,v$ در
مسیریابی تقاضای iام استفاده شده است یا خیر.

\begin{align}
\begin{split}
	min \sum_{i = [1, ..., |D|]}\sum_{(u,v) \in E} x_{i,(u,v)}W_{(u,v)}
\end{split}
\end{align}

\begin{align}
\begin{split}
	\sum_{(u,v) \in E} x_{i,(u,v)} - \sum_{(v,u) \in E} x_{i,(v,u)} = 
	\left \{
		\begin{array}{cc}
			1 & u = s_i\\
			0 & u \in V - {s_i, t_i}\\
			-1 & u = t_i
		\end{array}
	\right.
	\\
	\quad u \in V, i = {1, ..., |D|}
\end{split}
\end{align}

\begin{align}
\begin{split}
	\sum_{i = [1, ..., |D|]}\sum_{(u,v) \in E} x_{i,(u,v)}Z_{(u,v)} \le d_i
\end{split}
\end{align}

\subsection{ب}

\subsection{ج}
\begin{align}
\begin{split}
	min \sum_{i = [1, ..., |D|]}\sum_{(u,v) \in E} x_{i,(u,v)}W_{(u,v)} + \lambda^T \sum_{i = [1, ..., |D|]}(\sum_{(u,v) \in E} x_{i,(u,v)}Z_{(u,v)} - d_i)
\end{split}
\end{align}

\begin{align}
\begin{split}
	\sum_{(u,v) \in E} x_{i,(u,v)} - \sum_{(v,u) \in E} x_{i,(v,u)} = 
	\left \{
		\begin{array}{cc}
			1 & u = s_i\\
			0 & u \in V - {s_i, t_i}\\
			-1 & u = t_i
		\end{array}
	\right.
	\\
	\quad u \in V, i = {1, ..., |D|}
\end{split}
\end{align}

\section{سوال سوم}
\subsection{الف}
\paragraph{}
متغیر $x_{i,(u,v)}$ نشان می‌دهد که آیا کانال i به یال $(u,v)$
اختصاص یافته است یا خیر.

\paragraph{}
متغیر $y_{i,u}$ نشان می‌دهد
که آیا کانال i در راس u استفاده شده است یا خیر

\begin{align}
\begin{split}
	\sum_{i=1}^{12} x_{i,(u,v)} = 1
	\quad (u,v) \in E
\end{split}
\end{align}

\begin{align}
\begin{split}
	x_{i,(u,v)} \le y_{i,u}\\
	x_{i,(u,v)} \le y_{i,v}\\
	\quad (u,v) \in E, 1 \le i \le 12
\end{split}
\end{align}

\begin{align}
\begin{split}
	y_{i,u} \le r_u
	u \in V, 1 \le u \le 12
\end{split}
\end{align}




\section{سوال چهارم}
\subsection{الف}



\section{سوال پنچم}

\end{document}